\section{A sharp Bernoulli likelihood transfer}\label{sec:sharp}

We give the probability step explicitly to distinguish a risk factor from
an exponential-rate factor. Sharp Chernoff comparisons in community
detection are also developed by \citet{zhouli}.
The argument below allows a perturbation measurable from independent
training data.

\begin{lemma}[Bounded tilted error preserves the sharp prefactor]
\label{lem:sharp}
Let $X_j$ be independent Bernoulli variables with probabilities $p_j$.
Let $T$ be an affine function of $(X_j)_j$, with linear coefficients
$w_j$, and put $\psi(\alpha)=\log\E e^{-\alpha T}$.
Suppose $\abs{w_j}\le C$, the minimizer $\alpha^*$ satisfies
$c\le \alpha^*\le1-c$ for a fixed $c>0$, and under the probability
measure $\mathbb{Q}$ proportional to $e^{-\alpha^*T}$,
\begin{equation}\label{eq:sharp-premises}
 \E_{\mathbb{Q}} T=0,\quad c{np}\le\operatorname{Var}_{\mathbb{Q}}(T)\le C{np},
 \quad\sum_jr_j\le C{np},
\end{equation}
where $r_j=\mathbb{Q}(X_j=1)$. Define
\[
 I=-\psi(\alpha^*)+\frac12\log(np).
\]
Within this lemma, $I$ is the exact transform-based exponent of $T$.
The added logarithm incorporates the square-root prefactor; hence
$I-\tfrac12\log(np)$ is the negative log transform at the optimizing
tilt. Its comparison with the full-size $I_{k\ell}$ of
Section~\ref{sec:oracle-scale} is made in the application below.
For every fixed $L<\infty$ and all sufficiently large ${np}$, uniformly over
$|x|\le L$,
\begin{equation}\label{eq:sharp-oracle}
 e^{-I}\lesssim\PP(T<x)
 \le\PP(T\le x)\lesssim e^{-I}.
\end{equation}
The implicit constants may depend on $L$ and the fixed constants in
the oracle hypotheses, but not on $n$ or on $x$ within this interval.

Let $\mathcal W$ be training data independent of $(X_j)_j$ and let
\[
 \widehat T=T+\varepsilon_0+\sum_j\varepsilon_jX_j
\]
have coefficients measurable with respect to $\mathcal W$. On a measurable
training event suppose
\begin{equation}\label{eq:tilt-error}
 \max_j|\varepsilon_j|\le R,\qquad
 \left|\varepsilon_0+\sum_jr_j\varepsilon_j\right|\le B_1,\qquad
 \sum_jr_j\varepsilon_j^2\le B_2.
\end{equation}
For fixed $R,B_2$ and every real $x$,
\begin{equation}\label{eq:sharp-transfer}
 \PP(\widehat T\le x\mid\mathcal W)
 \lesssim\exp\{\alpha^*x+B_1+e^R B_2/2\}\,e^{-I}.
\end{equation}
The implicit factor may depend on the fixed oracle bounds and on
$R,B_2$, but is uniform in $x,B_1,n$ and over the training event.
In particular $B_1$ can
be a nonnegative training-measurable random bound, provided its indicated
exponential moment is integrated when removing the conditioning.
\end{lemma}

\begin{proof}
Under $\mathbb{Q}$, the coordinates remain independent Bernoulli variables,
\[
 r_j=\frac{p_je^{-\alpha^*w_j}}{1-p_j+p_je^{-\alpha^*w_j}}.
\]
The sum of their centered weighted third absolute moments is at most
$C\sum_jr_j=O({np})$. Berry--Esseen and
\eqref{eq:sharp-premises} therefore show that every interval of length one
has $\mathbb{Q}$ probability at most $C/\sqrt{np}$, uniformly in
its location. The exact change of measure is
\[
 \frac{\mathrm{d}\PP}{\mathrm{d}\mathbb{Q}}
 =\exp\{\psi(\alpha^*)+\alpha^*T\}
 =\sqrt{np}\,e^{-I+\alpha^*T}.
\]
Partitioning $(-\infty,x]$ into unit intervals cancels the
$\sqrt{np}$ factor against the interval-probability bound and gives
the upper bound in \eqref{eq:sharp-oracle}.

For the lower bound choose a fixed interval $(x-L_0,x)$, with $L_0$
sufficiently large. Its Gaussian probability under a centered normal
variable of variance in $[c{np},C{np}]$ is at least $cL_0/\sqrt {np}$ for large
${np}$, uniformly over bounded $x$. Choose $L_0$ so this exceeds twice the
Berry--Esseen error. The $\mathbb{Q}$ probability is then at least $c'/\sqrt{np}$.
On this interval $e^{\alpha^*T}$ has a fixed positive lower bound,
so the same change-of-measure identity gives $c_Le^{-I}$.
Using strict endpoints also handles arbitrary lattice atoms and tie rules.

For the perturbation, independence of training and the row allows us to
regard $\varepsilon_0,\varepsilon$ as fixed. The elementary Bernoulli exponential-moment bound
gives
\begin{align*}
 \log\E e^{-\alpha^*\widehat T}
 &= -I+\tfrac12\log(np)-\alpha^*\varepsilon_0+\log\E_{\mathbb{Q}}e^{-\alpha^*\sum_j\varepsilon_jX_j}\\
 &\le -I+\tfrac12\log(np)-\alpha^*\left(\varepsilon_0+\sum_jr_j\varepsilon_j\right)
       +\frac{(\alpha^*)^2e^{\alpha^*R}}2\sum_jr_j\varepsilon_j^2\\
 &\le -I+\tfrac12\log(np)+B_1+e^R B_2/2.
\end{align*}
Now tilt the original law by $e^{-\alpha^*\widehat T}$.
Write $\widehat r_j=\widehat{\mathbb Q}(X_j=1)$.
Under this new measure $\widehat{\mathbb{Q}}$ the coordinates remain independent and
their odds are those under $\mathbb{Q}$ multiplied by $e^{-\alpha^*\varepsilon_j}$.
Consequently their probabilities and variances are within fixed factors
of those under $\mathbb{Q}$. Furthermore,
\[
 \left(\sum_jr_j(1-r_j)(w_j+\varepsilon_j)^2\right)^{1/2}
 \ge \sqrt{c{np}}-\sqrt{B_2}.
\]
For large ${np}$, the variance of $\widehat T$ under $\widehat{\mathbb{Q}}$ is therefore
between positive constant multiples of ${np}$. Its centered third absolute
moments sum to $O({np})$ because $w_j+\varepsilon_j$ are uniformly bounded and
$\sum_j\widehat r_j=O({np})$. Berry--Esseen bounds every unit interval under
$\widehat{\mathbb{Q}}$ by $C/\sqrt {np}$, irrespective of its mean. Multiplication by the displayed exponential-moment bound
cancels the $1/\sqrt{np}$ interval factor. The same interval summation
therefore proves \eqref{eq:sharp-transfer}.
\end{proof}

\subsection{Application to the actual Bernoulli SBM oracle}

Fix a tested vertex with true label $k$ and realize the count vector
from Section~\ref{sec:statements} using its actual incident edges.
For this fixed target, $T_{k\ell}$ is therefore exactly its raw-row
contrast, not an independent copy. Its distribution depends only on
$k,\ell$, the community sizes and $P$. The vertex is still indexed by
$i$ wherever its deleted graph or projection is used.
For the projection arguments the linear coefficient vector is
\begin{equation}\label{eq:sbm-weights}
 (w_{k\ell})_j
   =\log\frac{P_{kz_j}(1-P_{\ell z_j})}
                    {P_{\ell z_j}(1-P_{kz_j})},\qquad j\in[n].
\end{equation}
Thus, for the chosen target,
\[
 T_{k\ell}=\sum_{j\ne i}A_{ij}(w_{k\ell})_j
       +\sum_{j\ne i}\log\frac{1-P_{kz_j}}{1-P_{\ell z_j}}.
\]
The nonedge term is retained, and $w_{k\ell,-i}$ means that coordinate
$i$ is deleted. The independent binomial counts give the exact transform
\begin{equation}\label{eq:sbm-information}
 \log\E_k e^{-\alpha T_{k\ell}}
 =\sum_{r=1}^K(n_{r}-\ind{r=k})\,
 \log\!\left[P_{kr}^{1-\alpha}P_{\ell r}^\alpha
 +(1-P_{kr})^{1-\alpha}(1-P_{\ell r})^\alpha\right].
\end{equation}
The self-loop contributes nothing. To compare this exact transform
with the full-size expression in Section~\ref{sec:oracle-scale}, note
that the additional affinity factor obeys, for $0\le\alpha\le1$,
\[
 1-p\le P_{kk}^{1-\alpha}P_{\ell k}^{\alpha}
 +(1-P_{kk})^{1-\alpha}(1-P_{\ell k})^{\alpha}\le1.
\]
The lower bound follows from the nonedge term, and the upper bound
from H\"older's inequality. Taking negative logarithms and maximizing
shows that
\[
 0\le I_{k\ell}
       +\min_{0\le\alpha\le1}\log\E_k e^{-\alpha T_{k\ell}}
       -\tfrac12\log(np)\le-\log(1-p).
\]
Thus the exact transform-based exponent in \cref{lem:sharp} is
$I_{k\ell}+O(p)$, and its exponential lies between $e^{-I_{k\ell}}$
and $(1-p)^{-1}e^{-I_{k\ell}}$. Under $np^2\le C_1$, eventually
$p\le1/2$, so this additional multiplier is uniformly at most two.
Every optimizing tilt and change-of-measure identity in the lemma
still uses the exact transform \eqref{eq:sbm-information}, not a
fictitious observed self-loop. The optimizer of the full-size
expression need not coincide with that exact optimizer.
On the compact positive full-rank classes used here, probabilities are
comparable to $p$ and bounded away from one, the weights are bounded,
$\alpha^*$ stays away from the endpoints, and the tilted variance is comparable
to ${np}$. To verify the last assertions, the derivatives of the transform
at its endpoints are respectively minus and plus the two directional
Bernoulli KL divergences, both of order ${np}$. Its second derivative is
bounded by $C{np}$. The row separation implied by
$\sigma_K(P)\ge c p$ gives a lower bound $c{np}$ on that derivative
throughout $[0,1]$: tilted probabilities are comparable to $p$ and the
sum of squared log-odds differences is at least $cn$.
Thus \cref{lem:sharp} applies uniformly to every pair. Combining it
with the one-coordinate comparison above gives, for every fixed
bounded threshold $x$, the bounds
\[
 \PP_k(T_{k\ell}\le x)\asymp e^{-I_{k\ell}},
 \qquad -\log\PP_k(T_{k\ell}\le x)=I_{k\ell}+O(1).
\]
These are constant-factor comparisons, not exact risk identities.
The use of $\tfrac12\log(np)$ is justified by the tilted variance
being comparable to $np$ and the optimizer staying away from the
endpoints on the stated parameter class. It is not a universal
prefactor formula for vanishing signal separation.

\subsection{Fixed-\texorpdfstring{$K$}{K} comparison and label alignment}

A strict pairwise defeat of the true community forces an oracle
error; an error entails a non-strict defeat by some competitor.
Under the joint community-$k$ count law this gives
\[
 \begin{split}
 \max_{\ell\ne k}\PP_k(T_{k\ell}<0)
 &\le\PP_k\!\left\{
       \min_{\ell<k}T_{k\ell}\le0
       \ \text{or}\ \min_{\ell>k}T_{k\ell}<0\right\}\\
 &\le\sum_{\ell\ne k}\PP_k(T_{k\ell}\le0).
 \end{split}
\]
The sharp strict and non-strict tails, together with fixed $K$, imply
\begin{equation}\label{eq:oracle-pairs}
 R_{\mathrm{oracle}}
 \asymp\sum_{k=1}^K\frac{n_k}{n}
                    \sum_{\ell\ne k}e^{-I_{k\ell}}.
\end{equation}
The weighting $n_k/n$ is the original average over vertices, regrouped
by their true community. This is not a redefinition as a worst-pair
oracle. The argument requires no independence among competing tests.
Fitted partitions are aligned with the true groups by a single common
permutation on the good event. Uniform score comparisons under that
alignment imply
\[
 \E[\Mis(\widehat z,z)\ind{\mathcal G}]
 \le\frac1n\sum_i\PP(\text{aligned fitted error at }i,\mathcal G).
\]
The ensuing pairwise bounds apply to fixed true coefficient vectors; a
data-dependent matching does not become an independence assumption.

\subsection{Exceptional coefficients}

\begin{lemma}[Paying exceptional coefficients only when the row hits them]
\label{lem:sparse-truncated-transfer}
Assume the oracle hypotheses of \cref{lem:sharp}, and write $p_j$
for the un-tilted row probabilities. Suppose $p_j\le Cr_j$, as
holds here for bounded oracle likelihood coefficients. On an independent
training event $\mathcal E$, suppose
\[
 \left|\varepsilon_0+\sum_jr_j\varepsilon_j\right|\le B_1,\qquad \sum_jr_j\varepsilon_j^2\le B_2.
\]
Here $B_1,B_2$ are fixed constants. For any fixed $R>0$ and bounded
$x$, the exceptional coordinates can be charged directly:
\begin{equation}\label{eq:sparse-truncated-transfer}
 \begin{split}
 \PP(\widehat T\le x)
 \le{}& C_Re^{-I}+\PP(\mathcal E^c)\\
 &+\E\!\left[\ind{\mathcal E}
             \sum_{j:|\varepsilon_j|>R}p_j\right].
 \end{split}
\end{equation}
The last term is the expected probability bound for the tested row
hitting an exceptional coefficient. Since
\[
 \sum_{j:|\varepsilon_j|>R}p_j
 \le\frac{C}{R^2}
      \sum_jr_j\varepsilon_j^2\ind{|\varepsilon_j|>R},
\]
a corresponding bound for the truncated quadratic error also suffices.
Either expected term of order $e^{-I}$ is enough.
If $I\le Cnp$ for a fixed $C$, an exceptional-term bound
$C_se^{-snp}$ is sufficient for any fixed $s>C$.
\end{lemma}

\begin{proof}
For analysis only, set
$\varepsilon'_j=\varepsilon_j\ind{|\varepsilon_j|\le R}$.
Then
\[
 \left|\varepsilon_0+\sum_jr_j\varepsilon'_j\right|
 \le B_1+\frac1R\sum_{j:|\varepsilon_j|>R}r_j\varepsilon_j^2
 \le B_1+B_2/R,\qquad
 \sum_jr_j(\varepsilon'_j)^2\le B_2.
\]
The bounded-coefficient lemma bounds the truncated score's error by
$C_Re^{-I}$ on $\mathcal E$.
The two random scores agree whenever $X_j=0$ on every exceptional
coordinate. Conditional on training, the probability they differ is
at most $\sum_{j:|\varepsilon_j|>R}p_j$. Integrating this bound and
paying $\PP(\mathcal E^c)$ proves the result. The truncation is a
proof device, not an additional operation in the decoder.
\end{proof}

